
\subsection{聚束最优传播、能散与退相干}\label{subsec:bunching-decoherence-closure-dp47}

\eqnrefs{eq:electron_bunching_harmonic_dp46} 已把散射后能量格密度矩阵映射为传播后的密度谐波。实验上真正需要的是在给定漂移距离、输入能散和环境噪声下的可见谐波。令 $t=L/v_0$，并把理想聚束因子记为 $B_{\Delta\ell}^{(0)}(L)$。若输入电子具有小的速度涨落 $\delta v$，则同一实验距离对应传播时间涨落
\begin{equation*}
 \delta t\simeq-\frac{L}{v_0^2}\delta v.
\end{equation*}
对窄带相对论电子，$\delta v=(\partial v/\partial E)_{E_0}\delta E$。定义
\begin{equation}
\boxed{
 \sigma_t(L)
 =\frac{L}{v_0^2}
 \left|\frac{\partial v}{\partial E}\right|_{E_0}\sigma_E .}
\label{eq:arrival_time_jitter_energy_spread_dp47}
\end{equation}
若第 $\Delta\ell$ 阶密度谐波的局部角频率为 $\Delta\ell\omega$，对高斯能量分布平均后得到
\begin{equation}
\boxed{
 B_{\Delta\ell}^{(E)}(L)
 =B_{\Delta\ell}^{(0)}(L)
 \exp\!\left[-\frac12\Delta\ell^2\omega^2\sigma_t^2(L)\right].}
\label{eq:bunching_energy_spread_damping_dp47}
\end{equation}
因此高阶谐波比基频更快被有限能散压低，其衰减指数严格按 $\Delta\ell^2$ 增长。

开放系统退相干可以在同一对象上表示。若能量格的 Markov 纯退相干满足
\begin{equation}
\boxed{
 \dot\rho_{\ell\ell'}\big|_{\rm deph}
 =-\Gamma_\varphi(\ell-\ell')^2\rho_{\ell\ell'}.}
\label{eq:bunching_pure_dephasing_master_dp68}
\end{equation}
则第 $\Delta\ell$ 条副对角线统一获得因子
\begin{equation}
\boxed{
 B_{\Delta\ell}^{(\mathrm{deph})}(L)
 =B_{\Delta\ell}^{(0)}(L)
 \exp\!\left[-\Delta\ell^2\Gamma_\varphi\frac{L}{v_0}\right].}
\label{eq:bunching_dephasing_damping_dp47}
\end{equation}
能散平均与环境退相干来自不同物理机制，但都作用于同一个 $B_{\Delta\ell}$。若二者独立，实验前向模型可写成
\begin{equation}
\boxed{
 B_{\Delta\ell}^{\rm obs}(L)
 =B_{\Delta\ell}^{(0)}(L)
 \exp\!\left[-\frac12\Delta\ell^2\omega^2\sigma_t^2(L)
 -\Delta\ell^2\Gamma_\varphi\frac{L}{v_0}\right].}
\label{eq:bunching_observed_master_dp47}
\end{equation}
这一定量式把电子源品质、漂移长度和环境相干时间直接连接到可观测密度谐波。

对给定调制强度，最优漂移长度应由最大化目标谐波而不是引用固定 Talbot 分数来确定。定义
\begin{equation}
\boxed{
 L_{\Delta\ell}^*\equiv\underset{L>0}{\operatorname{argmax}}
 \left|B_{\Delta\ell}^{\rm obs}(L)\right| .}
\label{eq:optimal_bunching_length_dp47}
\end{equation}
在无噪声二阶色散模型中，$B_{\Delta\ell}^{(0)}(L)$ 对 $L_T$ 周期；加入\eqnrefs{eq:bunching_observed_master_dp47} 的指数包络后，不同 Talbot 周期不再等价，通常应选择最早达到局域极大的漂移距离。

若最终实验是用整形电子驱动目标跃迁，则\secref{subsec:shaped-electron-drive-dp19}的关系直接给
\begin{equation}
\boxed{
 \Omega_{3i}^{(e)}(L)
 =\frac{2G_{3i}}{T}B_{\Delta\ell_i}^{\rm obs}(L),
 \qquad
 \omega_{3i}=\Delta\ell_\ii\omega .}
\label{eq:target_rabi_with_decoherence_dp47}
\end{equation}
于是电子源能散和电子退相干不再是独立的“束流质量指标”，而是直接进入目标 Rabi 频率的乘法衰减。


